\section{Introduction}

\subsection{The Problem in a Question}
The paper \textit{Efficient Tagged Memory} \cite{joannou2017efficient} describes how tagged memory architectures have been explored in academia and employed in industry many times with various different focuses throughout the history of computer design. Through both academic rigour and market testing they have been proven to provide system-wide security properties such as hardware capabilities \cite{clark1989application}, pointer integrity \cite{crandall2004minos}, watchpoints \cite{greathouse2012case}, and information-flow tracking \cite{suh2004secure}.\\
The employment of a tagged memory architecture can certainly provide utility, but there are associated costs that render widespread adoption of tagged memory architectures currently infeasible. With a simplistic scheme, the number of memory accesses that workloads must perform doubles due to the CPU now having to access the tag of the value it writes, as well as the value itself, which brings an unacceptably large time overhead. Also, the table in which tags are stored comes with a constant memory consumption cost, which is particularly disadvantageous within the server-side compute industry, where every byte of DRAM is precious.

Efforts have been and are currently being made to mitigate the overheads of tagged memory. The question that this dissertation is concerned with is:
\begin{center}
	``How can we further reduce the overheads that tagged memory introduces?''
\end{center}

\subsection{Potential Answers}
Tag caches are often used in practice to hold tag data close to the CPU for faster access times. The main area in which mitigation of the memory access latency overhead is found is in increasing the `cacheability' of the tag table in DRAM. That is, by increasing the amount of information about the tags that we store in a fast-access cache, we reduce the average latency overhead that is introduced by tagged memory's new requirement of accessing tag data, as well as value data. An effort of this kind, described in \textit{Efficient Tagged Memory}, has been shown to reduce the common-case DRAM traffic overhead to ``well below 1\%''.\\
Also, suppose that the tagged memory system is one which runs a hypervisor, which is a common setup for servers. Suppose further that the hypervisor has a dynamic DRAM allocator that could offer free physical address regions to its guest operating systems as and when they require it. A novel system involving a dynamically-resizing tag table that can be compressed and decompressed at runtime would allow the guest OSes to reserve and free regions of DRAM that they use for storing their tags as they execute, meaning that each OS will consume large amounts of extra space only when necessary, as opposed to all of them reserving at boot-time the maximum that they might require during any part of their execution. By increasing the `compressibility' of the tag table, we can keep the size of the tag table in each guest OS instance closer to a theoretical minimum, given the current workloads each is performing.

\subsection{My Project}
These potential areas of improvement motivate the exploration and investigation of existing tag controller schemes in search of opportunities to improve upon them, to increase the `cacheability' and `compressibility' of the DRAM tag table. They also motivate the creation of tools that allow architects to quantitatively and qualitatively evaluate new tag controller schemes that they might design.\\
My project aims to explore existing tag controller implementations, as well as provide a framework for an architect to implement a controller of their own design, which provides them with qualitative and quantitative data that allows them to usefully investigate and understand its performance, for the sake of contributing to the field of research in this area, and participating in answering the question I posed earlier.
